% ============ 30 主引理 (原书 pp.176-183) ============
\chapter{主引理}\label{sec:30}
\index{主引理}
\index{权}
\index{辛群}
\index{单值}
\index{不变理论}

\subsection*{记号回顾}

$U_0$ 又是 $\mathbb{F}_q$ 上的非奇异仿射曲线，$U$ 表示 $U_0$ 看作 $\mathbb{F}$ 上的曲线.于是对某个仿射 $\mathbb{F}_q$-代数 $A_0$，$U_0 = \operatorname{Specm} A_0$，我们设 $K$ 是 $A_0$ 的分式域.于是 $K$ 是 $U_0$ 上的有理函数域.选取 $K$ 的代数闭包 $\Omega$，把 $\pi_1(U_0)$ 等同于 $\operatorname{Gal}(K^{\mathrm{un}}/K)$($K^{\mathrm{un}}$ 是 $K$ 在 $\Omega$ 中的最大非分歧扩张(相对于 $A_0$))，并把 $\pi_1(U)$ 等同于 $\operatorname{Gal}(K^{\mathrm{un}}/K \cdot \mathbb{F})$.则 $U_0$ 上 $\mathbb{Q}_{\ell}$-向量空间的局部常值层 $E_0$ 对应于带 $\pi_1(U_0)$ 连续作用的有限维 $\mathbb{Q}_{\ell}$-向量空间 $E$.$E$ 在 $U$ 上的逆像对应于看作 $\pi_1(U)$-模的 $E$.这里 $E$ 是 $E$ 在由 $\operatorname{Specm} \Omega \to \operatorname{Specm} K \to U_0$ 定义的 $U_0$ 的几何点处的茎 $E_{\bar{\xi}}$.

有正合序列
\[
0 \to \pi_1(U) \to \pi_1(U_0) \to \operatorname{Gal}(\mathbb{F}/\mathbb{F}_q) \to 0.
\]
对每个闭点 $x \in U_0$，有 Frobenius 元素 $\varphi_x \in \pi_1(U_0)$(良定义到共轭)，它固定 $A$ 在 $K^{\mathrm{un}}$ 中的整闭包 $A^{\mathrm{un}}$ 的位于 $\mathfrak{p}$ 之上的某个素理想 $\mathfrak{P}$，并在 $A^{\mathrm{un}}/\mathfrak{P}$ 上按 $a \mapsto a^{q^{\deg x}}$ 作用.$\varphi_x$ 在 $\operatorname{Gal}(\mathbb{F}/\mathbb{F}_q)$ 中的像是 $\varphi^{\deg x}$.

沿用 Deligne，我设
\[
F_x = \varphi^{-1}_x \in \pi_1(U_0), \qquad F = \varphi^{-1} \in \operatorname{Gal}(\mathbb{F}/\mathbb{F}_q).
\]
于是 $F_x$ 与 $F$ 分别在 $E_x = E$ 与 $H^{r}_{c}(U, E)$ 上的作用与它们同名者几何定义的作用一致(29.10; 29.12).

\subsection*{线性代数预备}

\begin{lemma}\label{lem:30.1}
设 $\alpha$ 与 $\beta$ 分别是有限维 $K$-向量空间 $V$ 与 $W$ 的自同态.则 $\alpha \otimes \beta \colon V \otimes W \to V \otimes W$ 的迹是 $\alpha$ 与 $\beta$ 的迹的乘积.
\end{lemma}

\begin{proof}
分别选取 $V$ 与 $W$ 的基 $(e_i)$ 与 $(f_i)$，并设
\[
\alpha e_i = \sum a_{ji} e_j; \qquad \beta f_i = \sum b_{ji} f_j.
\]
则 $(e_i \otimes f_{i'})$ 是 $V \otimes W$ 的基，且
\[
(\alpha \otimes \beta)(e_i \otimes f_{i'}) = \sum_{j, j'} a_{ji} b_{j'i'} e_j \otimes f_{j'}.
\]
因此，$\alpha \otimes \beta$ 的迹是
\[
\operatorname{Tr}(\alpha \otimes \beta) = \sum_{i, i'} a_{ii} b_{i'i'} = \left( \sum a_{ii} \right) \left( \sum b_{ii} \right) = (\operatorname{Tr}(\alpha))(\operatorname{Tr}(\beta)).
\]
\end{proof}

\begin{lemma}\label{lem:30.2}
设 $K$ 是含 $\mathbb{Q}$ 的域，$\alpha \colon V \to V$ 是 $K$-向量空间 $V$ 的自同态.$\alpha$ 的特征多项式有 $\mathbb{Q}$ 中的系数，当且仅当存在一组基，使 $\alpha$ 的矩阵元素在 $\mathbb{Q}$ 中.
\end{lemma}

\begin{proof}
当然，若 $\alpha$ 关于 $V$ 的某组基的矩阵 $A$ 的元素在 $\mathbb{Q}$ 中，则特征多项式 $\det(1 - At)$ 在 $\mathbb{Q}[t]$ 中.逆命题由矩阵的有理标准形理论推出.
\end{proof}

\begin{definition}\label{def:30.3}
含 $\mathbb{Q}$ 的域上向量空间的自同态称为有理的，如果它满足引理~\ref{lem:30.2} 的等价条件.
\end{definition}

对向量空间 $V$，$V^{\otimes m}$ 表示 $V$ 的 $m$ 个副本的张量积.$V$ 的自同态 $\alpha$ 定义 $V^{\otimes m}$ 的自同态 $\alpha^{\otimes m}$.由引理~\ref{lem:30.1}，看到
\[
\operatorname{Tr}(\alpha^{\otimes m}|V^{\otimes m}) = (\operatorname{Tr}(\alpha|V))^{m}.
\]
显然 $\alpha$ 有理时 $\alpha^{\otimes m}$ 也有理.我们常把 $\alpha^{\otimes m}$ 写为 $\alpha$.

\subsection*{定理的陈述}

\begin{definition}\label{def:30.4}
设 $n \in \mathbb{Z}$.

(a) $U_0$ 上 $\mathbb{Q}_{\ell}$-向量空间的局部常值层 $E_0$ 称为有权 $n$，如果对 $U_0$ 的所有闭点 $x$，$F_x \colon E_x \to E_x$ 的每个特征值都是代数数，其复共轭有绝对值 $q^{\deg(x)n/2}$.

(b) 带 $\pi_1(U_0)$ 连续作用的有限维 $\mathbb{Q}_{\ell}$-向量空间 $E$ 有权 $n$，如果对 $U_0$ 的所有闭点 $x$(即 $A_0$ 的非零素理想 $\mathfrak{p}$)，$\varphi_x \colon E_x \to E_x$ 的每个特征值都是代数数，其复共轭有绝对值 $q^{-\deg(x)n/2}$.

当 $E$ 是对应于 $E_0$ 的 $\pi_1(U_0)$-模时，这些概念一致：$E_0$ 有权 $n$，当且仅当 $E$ 有权 $n$.
\end{definition}

\begin{example}\label{ex:30.5}
(a) 层 $\mathbb{Q}_{\ell}(1)$ 对应于 $\pi_1(U_0)$-模 $\left( \varprojlim_{n} \mu_{\ell^{n}}(K^{\mathrm{un}}) \right) \otimes_{\mathbb{Z}_{\ell}} \mathbb{Q}_{\ell}$.因为 $\ell \neq p$，$\mu_{\ell^{n}}(\mathbb{F}) = \mu_{\ell^{n}}(K^{\mathrm{un}})$ 是秩 1 的自由 $\mathbb{Z}/\ell^{n}\mathbb{Z}$-模：$m \in \mathbb{Z}/\ell^{n}\mathbb{Z}$ 按规则 $(m, \zeta) \mapsto \zeta^{m}$ 作用.因此 $\varprojlim_{n} \mu_{\ell^{n}}(\mathbb{F})$ 是秩 1 的自由 $\mathbb{Z}_{\ell}$-模，$\mathbb{Q}_{\ell}(1)$ 是维数 1 的 $\mathbb{Q}_{\ell}$-向量空间.群 $\pi_1(U_0)$ 通过它的商 $\operatorname{Gal}(\mathbb{F}/\mathbb{F}_q)$ 作用于 $\mathbb{Q}_{\ell}(1)$，$\operatorname{Gal}(\mathbb{F}/\mathbb{F}_q)$ 的典范生成元 $\varphi = (z \mapsto z^{q})$ 在 $\mathbb{Q}_{\ell}(1)$ 上按乘以 $q$ 作用.Frobenius 元素 $\varphi_x$ 在 $\operatorname{Gal}(\mathbb{F}/\mathbb{F}_q)$ 中的像是 $\varphi^{\deg x}$，从而在 $\mathbb{Q}_{\ell}(1)$ 上按 $q^{\deg x}$ 作用.于是，$\mathbb{Q}_{\ell}(1)$ 有权 $-2$.

(b) 若 $E_0$ 有权 $\beta$，则 $E^{\otimes m}_0$ 有权 $m\beta$：例如，$\mathbb{Q}_{\ell}(m)$ 有权 $-2m$.
\end{example}

下面的定理常称为主引理(限制形式).

\begin{theorem}\label{thm:30.6}
设 $E$ 是对应于 $\mathbb{Q}_{\ell}$-向量空间的局部常值层 $E_0$ 的 $\pi_1(U_0)$-模.设 $n$ 是整数.假定：
\begin{enumerate}
  \item (有理性.) 对所有闭点 $x \in U_0$，$F_x$ 在 $E$ 上的作用有理(\ref{def:30.3})；
  \item 存在非退化 $\pi_1(U_0)$-不变斜对称形式
  \[
  \psi \colon E \times E \to \mathbb{Q}_{\ell}(-n);
  \]
  \item (大几何单值.) $\pi_1(U)$ 在 $\operatorname{Sp}(E, \psi)$ 中的像(对 $\ell$ 进拓扑)开.
\end{enumerate}
则：
\begin{enumerate}
  \item $E$ 有权 $n$，即 $F_x$ 在 $E_x$ 上的特征值有绝对值 $(q^{\deg x})^{n/2}$.
  \item $F$ 在 $H^{1}_{c}(U, E)$ 上的作用有理，且它的特征值都有绝对值 $\leq q^{n/2+1}$.
  \item 设 $j$ 是 $U$ 到 $\mathbb{P}^{1}$ 的包含.$F$ 在 $H^{1}(\mathbb{P}^{1}, j_{*}E)$ 上的作用有理，且它的特征值 $\alpha$ 满足
  \[
  q^{n/2} < |\alpha| < q^{n/2+1}.
  \]
\end{enumerate}
这里 $\operatorname{Sp}(E, \psi)$ 表示 $\psi$ 的辛群，即
\[
\operatorname{Sp}(E, \psi) = \{\alpha \in \mathrm{GL}(E) \mid \psi(\alpha e, \alpha e') = \psi(e, e'), \ e, e' \in E\}.
\]
设 $\gamma \in \pi_1(U_0)$ 满足 $\gamma|_{\mathbb{F}} = \varphi^{m}$($m \in \mathbb{Z}$ 适当).$\psi$ 是 $\pi_1(U_0)$-不变的意思是
\[
\psi(\gamma e, \gamma e') = q^{-nm}\psi(e, e'), \qquad e, e' \in E.
\]
特别地，若 $\gamma \in \pi_1(U)$，则它作为辛群的元素作用于 $E$，因此 (c) 有意义.
\end{theorem}

\begin{example}\label{ex:30.7}
设 $d$ 是奇数，假定有正则映射
\[
\varphi \colon Y_0 \to U_0,
\]
使得对 $U_0$ 的每个闭点 $x$，$\varphi^{-1}(x)$ 是定义在 $\kappa(x) = \mathbb{F}_{q^{\deg x}}$ 上的 $\mathbb{P}^{d+1}$ 中的非奇异超曲面 $Y(x)_0$.设 $E = R^{d}\varphi_{*}\mathbb{Q}_{\ell}$.则 $E$ 局部常值(由紧合-光滑基变换定理~\ref{thm:20.2})，且对任何闭点 $x \in U_0$，$E_x \cong H^{d}(Y(x), \mathbb{Q}_{\ell})$(紧合基变换定理~\ref{thm:17.2}).此外 $F_x$ 在 $E_x$ 上按 $Y(x)$ 的 Frobenius 映射在 $H^{d}(Y(x), \mathbb{Q}_{\ell})$ 上的作用作用，因此
\[
\det(1 - F_x t|E_x) = \det(1 - Ft|H^{d}(Y(x), \mathbb{Q}_{\ell})).
\]
现在，对奇数 $r \neq d$，$H^{r}(Y(x), \mathbb{Q}_{\ell})$ 为零，对偶数 $r$ 等于 $H^{r}(\mathbb{P}^{d}, \mathbb{Q}_{\ell}) \cong \mathbb{Q}_{\ell}(-\frac{r}{2})$(见例~\ref{ex:16.4})，因此
\[
Z(Y(x)_0, t) = \frac{\det(1 - Ft|H^{d}(Y(x), \mathbb{Q}_{\ell}))}{(1 - t)(1 - qt) \cdots (1 - q^{d}t)}.
\]
因为 $Z(Y(x)_0, t) \in \mathbb{Q}(t)$(见推论~\ref{cor:27.8})，$\det(1 - Ft|H^{d}(Y(x), \mathbb{Q}_{\ell}))$ 有有理系数，从而 $E$ 满足定理的条件 (a).

有层的典范配对
\[
R^{d}\varphi_{*}\mathbb{Q}_{\ell} \times R^{d}\varphi_{*}\mathbb{Q}_{\ell} \to R^{2d}\varphi_{*}\mathbb{Q}_{\ell} \cong \mathbb{Q}_{\ell}(-d),
\]
它在每个茎上变成杯积
\[
H^{d}(Y(x), \mathbb{Q}_{\ell}) \times H^{d}(Y(x), \mathbb{Q}_{\ell}) \to H^{2d}(Y(x), \mathbb{Q}_{\ell}) \cong \mathbb{Q}_{\ell}(-d).
\]
因此，层上的配对对应于配对
\[
\psi \colon E \times E \to \mathbb{Q}_{\ell}(-d),
\]
它是斜对称的(因为 $d$ 奇数)，非退化的(由 $\varphi$ 的由 $\operatorname{Spec} \Omega \to U_0$ 定义的几何一般纤维的 Poincar\'e 对偶)且 $\pi_1(U_0)$-不变的(因为它由 $U_0$ 上层 的态射定义).因此，$E$ 满足定理的条件 (b).若它满足条件 (c)，则它有权 $d$，且对每个 $Y(x)$，Riemann 猜想成立！

事实上，可以实现 $\mathbb{F}_q$ 上每个奇维数非奇异超曲面 $Y_0$ 为满足 (c) 的族(在可能扩张 $\mathbb{F}_q$ 后)的成员，从而这样证明 $Y_0$ 的 Riemann 猜想.这是几何问题：若我们能找到满足 (c) 的，含 $Y_0/\mathbb{F}$ 为成员的 $\mathbb{F}$ 上非奇异超曲面的族 $\varphi \colon Y \to U \subset \mathbb{A}^{1}$，则这个族将定义在 $\mathbb{F}_q$ 的有限扩张上.

设 $\delta$ 是 $Y_0$ 的次数.我们考虑 $d + 2$ 个变量的所有次数为 $\delta$ 的齐次多项式(到乘以非零标量)的集合.这个集合可以等同于 $\mathbb{P}^{N}$, $N = \binom{d+\delta+1}{\delta}$(见 AG 6.20).我们得到图
\[
\begin{tikzcd}
H \arrow[r, hook] \arrow[d, "\pi"'] & \mathbb{P}^{N} \times \mathbb{P}^{d+1} \arrow[ld] \\
\mathbb{P}^{N}
\end{tikzcd}
\]
使得对每个 $P \in \mathbb{P}^{N}$，$\pi^{-1}(P)$ 是由 $P$(看作齐次多项式)定义的 $\mathbb{P}^{d+1}$ 中的超曲面(可能是可约的).现在 $Y_0/\mathbb{F}$ 对应于 $\mathbb{P}^{N}$ 中的点 $P$，问题是证明存在过 $P$ 的直线，使``几何单值''大.
\end{example}

\subsection*{定理 (a) 的证明概述}

回忆
\[
Z(U_0, E^{\otimes 2k}, t) \overset{\mathrm{def}}{=} \prod_{x \in U_0} \frac{1}{\det(1 - F_x t^{\deg x}|E^{\otimes 2k})} = \prod_{x} \left( \sum_{m} \operatorname{Tr}(F^{m}_x|E^{\otimes 2k}) \frac{t^{m}}{m} \right).
\]
第一个等式是定义，第二个由初等引理~\ref{lem:27.5} 推出.定理的条件 (a) 蕴含 $Z(U_0, E^{\otimes 2k}, t) \in \mathbb{Q}[[t]]$，因此谈论它对 $t \in \mathbb{C}$ 的收敛半径有意义.

我们将(在定理的假设下)证明：
\begin{quote}
(I) 对所有正整数 $k$，$(\Omega^{2k}E)^{\pi_1(U)}$ 同构于 $\mathbb{Q}_{\ell}(-kn)$ 的若干副本的直和.

(A) 若对所有正整数 $k$，$Z(U_0, E^{\otimes 2k}, t)$ 对 $|t| < \frac{1}{q^{kn+1}}$ 收敛，则 $E$ 有权 $n$.
\end{quote}
我们解释这两个论断如何蕴含定理的 (a) 部分.如已经注意到的，$Z(U_0, E^{\otimes 2k}, t)$ 是系数在 $\mathbb{Q}$ 中的幂级数.Lefschetz 不动点公式的乘法形式~(\ref{thm:29.6})，
\[
Z(U_0, E^{\otimes 2k}, t) = \frac{\det(1 - Ft|H^{1}_{c}(U, E^{\otimes 2k}))}{\det(1 - Ft|H^{2}_{c}(U, E^{\otimes 2k}))},
\]
表明 $Z(U_0, E, t) \in \mathbb{Q}_{\ell}(t) \cap \mathbb{Q}[[t]] = \mathbb{Q}(t)$.回忆命题~\ref{prop:29.9}: $H^{2}_{c}(U, E) \cong E_{\pi_1(U)}(-1)$，因此 (I) 蕴含
\[
\det(1 - Ft|H^{2}_{c}(U, E^{\otimes 2k})) = (1 - q^{kn+1}t)^{N}, \qquad \text{某个 } N.
\]
于是
\[
Z(U_0, E^{\otimes 2k}, t) = \frac{\mathbb{Q}[t] \text{ 中的多项式}}{(1 - q^{kn+1}t)^{N}}.
\]
显然，它对 $|q^{kn+1}t| < 1$，即 $|t| < 1/q^{kn+1}$ 收敛.现在可以应用 (A) 表明 $E$ 有权 $n$.

\subsection*{(A) 的证明}

本节中，$E$ 满足定理的条件 (a) 与 (b).

\begin{lemma}\label{lem:30.8}
对所有正整数 $k$，幂级数
\[
\frac{1}{\det(1 - F_x t^{\deg x}|E^{\otimes 2k})}
\]
的系数是正有理数.
\end{lemma}

\begin{proof}
按初等引理~\ref{lem:27.5}，
\[
\log \frac{1}{\det(1 - F_x t^{\deg x}|E^{\otimes 2k})} = \sum_{m=1}^{\infty} \operatorname{Tr}(F^{m}_x|E^{\otimes 2k}) \frac{t^{m \deg(x)}}{m}.
\]
但 $\operatorname{Tr}(F^{m}_x|E^{\otimes 2k}) = \operatorname{Tr}(F^{m}_x|E)^{2k}$.在定理的条件 (a) 下，$\operatorname{Tr}(F^{m}_x|E)^{2} \in \mathbb{Q}$，因此 $\operatorname{Tr}(F^{m}_x|E^{\otimes 2k})$ 是正有理数.于是幂级数 $\log \frac{1}{\det(1 - F_x t^{\deg x}|E^{\otimes 2k})}$ 的系数是正有理数，它的 $\exp$ 的系数也是.
\end{proof}

\begin{lemma}\label{lem:30.9}
若对所有正整数 $k$，$Z(U_0, E^{\otimes 2k}, t)$ 对 $|t| < \frac{1}{q^{kn+1}}$ 收敛，则 $E$ 有权 $n$.
\end{lemma}

\begin{proof}
考虑
\[
Z(U_0, E^{\otimes 2k}, t) \overset{\mathrm{def}}{=} \prod_{x \in U_0} \frac{1}{\det(1 - F_x t^{\deg x}|E^{\otimes 2k})}.
\]
若 $a_m$ 是 $Z(U_0, E^{\otimes 2k}, t)$ 的幂级数展开中 $t^{m}$ 的系数，$a_{m, x}$ 是 $1/\det(1 - F_x t^{\deg x}|E^{\otimes 2k})$ 的展开中 $t^{m}$ 的系数，则 $a_m \geq a_{m, x}$(因为 $a_{0, x} = 1$ 且对所有 $x$，$a_{m, x} \geq 0$).因此 $Z(U_0, E^{\otimes 2k}, t)$ 的收敛半径 $\leq 1/\det(1 - F_x t^{\deg x}|E^{\otimes 2k})$ 的收敛半径.引理的假设因此蕴含 $1/\det(1 - F_x t^{\deg x}|E^{\otimes 2k})$ 对 $|t| < 1/q^{kn+1}$ 收敛.

若 $\alpha$ 是 $F_x$ 在 $E$ 上的特征值，则 $\alpha^{2k}$ 是 $F_x$ 在 $E^{\otimes 2k}$ 上的特征值，因此 $\frac{1}{\alpha^{2k\deg x}}$ 是 $1/\det(1 - F_x t^{\deg x}|E^{\otimes 2k})$ 的极点.因此，
\[
\left| \frac{1}{\alpha^{2k\deg x}} \right| \geq \frac{1}{q^{kn+1}}
\]
从而 $|\alpha^{2k}| \leq (q^{\deg x})^{kn+1}$.取 $2k$ 次根并让 $k \to \infty$，发现对 $F_x$ 在 $E$ 上的所有特征值 $\alpha$, $|\alpha| \leq (q^{\deg x})^{n/2}$.定理 (b) 中配对的存在性表明，对每个特征值 $\alpha$，存在特征值 $\alpha'$，使得 $\alpha\alpha' = (q^{\deg x})^{n}$，这完成了证明.
\end{proof}

\subsection*{(I) 的证明}

\begin{lemma}\label{lem:30.10}
定理的条件 (c) 蕴含 $\pi_1(U)$ 在 $\operatorname{Sp}(\psi)$ 中 Zariski 稠密，从而
\[
\operatorname{Hom}(E^{\otimes 2k}, \mathbb{Q}_{\ell})^{\pi_1(U)} = \operatorname{Hom}(E^{\otimes 2k}, \mathbb{Q}_{\ell})^{\operatorname{Sp}(\psi)}.
\]
\end{lemma}

\begin{proof}
设 $\bar{\gamma}$ 是 $\pi_1(U)$ 在 $\operatorname{Sp}(\psi, \mathbb{Q}_{\ell})$ 中的像.因为 $\pi_1(U)$ 紧(它是 Galois 群)，$\bar{\gamma}$ 紧，从而在 $\operatorname{Sp}(\psi, \mathbb{Q}_{\ell})$ 中闭.Cartan 定理的 $\ell$ 进版本(Serre, \emph{Lie Algebras and Lie Groups}, LG 5.42)于是说它是 $\operatorname{Sp}(\psi, \mathbb{Q}_{\ell})$ 的 Lie 子群.因为它开，它作为 Lie 群与 $\operatorname{Sp}(\psi, \mathbb{Q}_{\ell})$ 有相同的维数.设 $G \subset \operatorname{Sp}$ 是 $\bar{\gamma}$ 在 $\operatorname{Sp}$ 中的 Zariski 闭包(于是 $G$ 是定义在 $\mathbb{Q}_{\ell}$ 上的，满足 $G(\mathbb{Q}_{\ell}) \supset \bar{\gamma}$ 的 $\operatorname{Sp}$ 的最小子簇).则 $G$ 是 $\operatorname{Sp}$ 的代数子群，且
\[
\operatorname{Tgt}_e(\bar{\gamma}) \subset \operatorname{Tgt}_e(G) \subset \operatorname{Tgt}_e(\operatorname{Sp}).
\]
但 $\operatorname{Tgt}_e(\bar{\gamma}) = \operatorname{Tgt}_e(\operatorname{Sp})$，因此 $\dim G = \dim \operatorname{Sp}$.因为 $\operatorname{Sp}$ 连通，这蕴含 $G = \operatorname{Sp}$，从而 $\bar{\gamma}$ 在 $\operatorname{Sp}$ 中 Zariski 稠密.

最后，设 $f \colon E^{\otimes 2k} \to \mathbb{Q}_{\ell}$ 是被 $\pi_1(U)$ 固定的线性映射.对 $g \in \operatorname{Sp}(\psi)$，$g$ 固定 $f$ 是关于 $g$ 的代数条件：若 $f$ 被 $\bar{\gamma}$ 固定，则它被 $\bar{\gamma}$ 的 Zariski 闭包固定.
\end{proof}

显然，
\[
\operatorname{Hom}(E^{\otimes 2k}, \mathbb{Q}_{\ell})^{\operatorname{Sp}(\psi)} = \operatorname{Hom}((E^{\otimes 2k})^{\operatorname{Sp}(\psi)}, \mathbb{Q}_{\ell}).
\]
若 $f_1, f_2, \ldots, f_N$ 是 $\operatorname{Hom}(E^{\otimes 2k}, \mathbb{Q}_{\ell})^{\operatorname{Sp}(\psi)}$ 的一组基，则映射
\[
a \mapsto (f_1(a), \ldots, f_N(a)) \colon E^{\otimes 2k} \to \mathbb{Q}^{N}_{\ell}
\]
诱导同构
\[
(E^{\otimes 2k})^{\operatorname{Sp}(\psi)} \to (\mathbb{Q}_{\ell})^{N}.
\]
我们将用不变理论选取一组基 $(f_i)$，使我们能看出 $\operatorname{Gal}(\mathbb{F}/\mathbb{F}_q)$ 的作用如何转移到 $(\mathbb{Q}_{\ell})^{N}$ 上的作用.

考虑如下一般问题：给定特征零域 $K$ 上的向量空间 $V$ 与 $V$ 上的非退化斜对称形式 $\psi$，什么是 $\operatorname{Sp}$-不变的线性形式 $f \colon V^{\otimes 2k} \to K$？即什么是 $\operatorname{Hom}(V^{\otimes 2k}, K)^{\operatorname{Sp}}$？注意线性形式 $f \colon V^{\otimes 2k} \to K$ 可以看作多重线性形式
\[
f \colon V \times V \times \cdots \times V \to K.
\]
例如，
\[
(v_1, \ldots, v_{2k}) \mapsto \psi(v_1, v_2) \cdots \psi(v_{2k-1}, v_{2k})
\]
是这样的多重线性形式，它显然 $\operatorname{Sp}$-不变.更一般地，对 $\{1, \ldots, 2k\}$ 的任意划分为 $k$ 个两两不交的集 $P \colon \{\{a_1, b_1\}, \ldots, \{a_k, b_k\}\}$($a_i < b_i$)，得到不变形式：
\[
f_P \colon V^{2k} \to K, \qquad (v_1, \ldots, v_{2k}) \mapsto \prod \psi(v_{a_i}, v_{b_i}).
\]

\begin{proposition}\label{prop:30.11}
不变形式 $f_P$ 张成 $\operatorname{Hom}(V^{\otimes 2k}, K)^{\operatorname{Sp}}$.
\end{proposition}

\begin{proof}
见 H. Weyl, \emph{The Classical Groups}, PUP, VI.1, 1939(修订版 1946).或者见 Fulton, W., and Harris, J., \emph{Representation Theory}, Springer Verlag, 1991，附录 F.
\end{proof}

$\operatorname{Hom}(V^{\otimes 2k}, K)^{\operatorname{Sp}}$ 的 $f_P$ 型基给出 $\pi_1(U_0)$-等变映射
\[
E^{\otimes 2k} \to \mathbb{Q}_{\ell}(-kn)^{N},
\]
诱导同构 $(E^{\otimes 2k})^{\pi_1(U)} \to \mathbb{Q}_{\ell}(-kn)^{N}$.这证明了 (I).

\subsection*{定理 (b) 的证明}

因为 $U$ 仿射，$H^{0}_{c}(U, E) = 0$.另一方面，
\[
H^{2}_{c}(U, E) = E_{\pi_1(U)}(-1),
\]
它为零，因为条件 (c) 蕴含 $\pi_1(U)$ 在 $\operatorname{Sp}(E, \psi)$ 中 Zariski 稠密，且没有 $\operatorname{Sp}(E, \psi)$-不变的线性形式 $E \to \mathbb{Q}_{\ell}$(或对任何奇数 $n$，在 $E^{\otimes n}$ 上也没有，由不变理论).因此，Lefschetz 不动点公式~\ref{thm:29.6} 表明
\[
Z(U_0, E_0, t) = \det(1 - Ft|H^{1}_{c}(U, E)).
\]
按定义
\[
Z(U_0, E_0, t) = \prod_{x \in U_0} \frac{1}{(1 - F_x t^{\deg x}|E)},
\]
因此
\[
\frac{1}{\det(1 - Ft|H^{1}_{c}(U, E))} = \prod_{x} \frac{1}{\det(1 - F_x t^{\deg x}|E)}.
\]
条件 (a) 蕴含右边展开时是 $t$ 的系数在 $\mathbb{Q}$ 中的幂级数.为完成证明，我们证明它对 $|t| < 1/q^{\frac{n}{2}+1}$ 收敛.

回忆复分析中，若 $p_1, p_2, \ldots$ 是都不为零的复数序列，则若部分乘积收敛到非零复数，称 $\prod p_n$ 收敛.无限乘积 $\prod (1 + a_n)$ 绝对收敛，当且仅当级数 $\sum |a_n|$ 收敛.

设 $d$ 是 $E$ 的维数，$a_{x, i}$($1 \leq i \leq d$)是 $F_x$ 在 $E$ 上的特征值，使得
\[
\frac{1}{Z(U_0, E_0, t)} = \prod_{x, i} (1 - a_{x, i} t^{\deg x}).
\]
我们将证明
\[
\sum_{x, i} |a_{x, i} t^{\deg x}|
\]
对 $|t| < 1/q^{\frac{n}{2}+1}$ 收敛.需要的两个事实是：
\begin{itemize}
  \item $|a_{x, i}| = (q^{\deg x})^{\frac{n}{2}}$；
  \item $U_0$ 上次数为 $m$ 的闭点 $x$ 的个数 $\leq q^{m}$.
\end{itemize}
第一个论断在 (a) 中证明.第二个由如下事实推出：$U_0$ 上每个次数为 $m$ 的闭点 $x$ 至少贡献 1(事实上恰好 $m$)个 $U_0(\mathbb{F}_{q^m})$ 的元素，而后者有 $\leq q^{m}$ 个元素，因为 $U_0 \subset \mathbb{A}^{1}$.取 $|t| = \frac{1}{q^{\frac{n}{2}+1+\varepsilon}}$($\varepsilon > 0$).第一个事实蕴含
\[
\sum_{i} |\alpha_{i, x} t^{\deg x}| \leq \frac{d}{(q^{\deg x})^{1+\varepsilon}},
\]
第二个事实蕴含
\[
\sum_{x} \frac{1}{(q^{\deg x})^{1+\varepsilon}} \leq \sum_{m} \frac{q^{m}}{q^{m(1+\varepsilon)}} = \sum_{m} \frac{1}{q^{m\varepsilon}} < \infty.
\]
注意证明只用到了 $E$ 有权 $n$ 且 $E^{\pi_1(U)} = 0$.

\subsection*{定理 (c) 的证明}

由序列
\[
0 \to j_{!}E \to j_{*}E \to i_{*}i^{*}j_{*}E \to 0
\]
(见命题~\ref{prop:8.15})的上同调序列，得到满射
\[
H^{1}_{c}(U, E) \to H^{1}(\mathbb{P}^{1}, j_{*}E).
\]
因此 (b) 蕴含 $F$ 在 $H^{1}(\mathbb{P}^{1}, j_{*}E)$ 上有理作用，且它的特征值满足
\[
|\alpha| < q^{\frac{n}{2}+1}.
\]
层 $\mathcal{L}E(1)$ 与 $E$ 满足同样的假设，只把 $n$ 换成 $-2 - n$.因此，$F$ 在 $H^{1}(\mathbb{P}^{1}, j_{*}\mathcal{L}E(1))$ 上有理作用，且它的特征值 $\beta$ 满足
\[
|\beta| < q^{\frac{-2-n}{2}+1} = q^{-\frac{n}{2}}.
\]
现在对偶定理(下面的 32.3)给出典范非退化配对
\[
H^{r}(\mathbb{P}^{1}, j_{*}E) \times H^{2-r}(\mathbb{P}^{1}, j_{*}\mathcal{L}E(1)) \to H^{2}(\mathbb{P}^{1}, \mathbb{Q}_{\ell}(1)) \cong \mathbb{Q}_{\ell}.
\]
因此，每个 $\alpha$ 是某个 $\beta$ 的倒数，从而
\[
|\alpha| > q^{\frac{n}{2}}.
\]
